\documentclass{article}

\usepackage[utf8]{inputenc}
\usepackage[T1]{fontenc}
\usepackage{amsmath,amssymb,amsthm}
\usepackage{mathtools}
\usepackage{hyperref}
\usepackage{booktabs}

\makeatletter
\newcommand{\citep}[2][]{%
  \begingroup
  \def\@tmpkey{#2}%
  \ifx\@tmpkey\@empty\else
    (\ifx\@empty#1\@empty\else #1, \fi
     \@nameuse{bibkey@\@tmpkey})%
  \fi
  \endgroup
}
\expandafter\def\csname bibkey@lasser2026micro\endcsname{Microfoundation}
\expandafter\def\csname bibkey@lasser2026paper8b\endcsname{Need~Sufficiency}
\expandafter\def\csname bibkey@lasser2026paper8a\endcsname{Revealed~Sacrifice}
\expandafter\def\csname bibkey@manheim2019\endcsname{Manheim~\&~Garrabrant~2019}
\expandafter\def\csname bibkey@elmhamdi2024\endcsname{El-Mhamdi~\&~Hoang~2024}
\expandafter\def\csname bibkey@majka2025\endcsname{Majka~\&~El-Mhamdi~2025}
\makeatother

\newtheorem{definition}{Definition}
\newtheorem{lemma}{Lemma}
\newtheorem{theorem}{Theorem}
\newtheorem{remark}{Remark}
\newtheorem{assumption}{Assumption}
\newtheorem{proposition}{Proposition}
\newtheorem{corollary}{Corollary}

\newcommand{\HHI}{\mathrm{HHI}}
\newcommand{\Pproxy}{P^{\mathrm{proxy}}}
\newcommand{\Ttruth}{T^{\mathrm{truth}}}
\newcommand{\epsnonres}{\varepsilon_{\mathrm{nonres}}}
\newcommand{\Lip}{\mathrm{Lip}}
\newcommand{\Heff}{H_{\mathrm{eff}}}
\newcommand{\Neff}{N_{\mathrm{eff}}}
\newcommand{\Var}{\mathrm{Var}}

\title{Concentration-Gap selection theorem, v1: \\
algebraic weighted-selection kernel + counterparty interpretation \\
+ Route-C transfer to Goodhart slack}
\author{Working draft for codex review}
\date{}

\begin{document}

\maketitle

\section{Goal}

The Microfoundation paper~\citep{lasser2026micro} states the
\textbf{Concentration-Gap Conjecture}: optimization pressure on
the proxy correlates with proxy-truth gap exploitation.
Microfoundation defers a formal proof.  This draft attempts to
discharge a \emph{scoped version} of the conjecture by:

\begin{enumerate}
\item Proving an \textbf{algebraic weighted-selection theorem}
  (Theorem~\ref{thm:algebraic-selection}): for any
  Hilbert-space-valued ``deviation field'' over an indexed
  population, the weighted aggregate's distortion from the
  population mean is upper-bounded above and lower-bounded below
  by the weights' Herfindahl index $\HHI$, under explicit
  structural assumptions.

\item Interpreting the algebraic theorem as a
  \textbf{counterparty-selection theorem} in trade-flow terms
  (\S\ref{sec:counterparty}), where the deviation field is
  $\Delta_c = U_c - W$ (counterparty utility minus the
  welfare-relevant truth).

\item Transferring the selection theorem's distortion bound into
  the \textbf{GFM Goodhart slack} via Microfoundation's Lipschitz
  transfer (\S\ref{sec:routeC}): the proxy-truth gap
  $\epsnonres$ is bounded above and below by $\HHI$-derived
  quantities under the structural conditions.
\end{enumerate}

The scoped theorem is \emph{not} the universal Concentration-Gap
Conjecture: it depends on three explicit structural conditions
(representativeness, bounded dispersion, non-collusion).  But
each of these is operationally auditable (with proxies the
deployment-tooling specification can verify), unlike the
opaque dynamics conjecture they replace.

\section{The algebraic weighted-selection kernel}

\subsection{Setup}

Let $\mathcal{V}$ be a real Hilbert space with norm $\|\cdot\|$
(in applications, $\mathcal{V}$ is the space of capability-poset
measures or its tangent space).  Let $\mathcal{C}$ be a
finite-or-countable index set (in applications, the set of
counterparties).  Let $\Delta : \mathcal{C} \to \mathcal{V}$ be a
\textbf{deviation field}, $c \mapsto \Delta_c$.

A \textbf{weighting} is a probability measure on $\mathcal{C}$:
$w \in \Delta(\mathcal{C})$ with $w_c \geq 0$ and $\sum_c w_c = 1$.

The \textbf{weighted aggregate distortion} is
\begin{equation}
\label{eq:aggregate-distortion}
  \Delta(w) \;:=\; \sum_{c \in \mathcal{C}} w_c \cdot \Delta_c
  \;\in\; \mathcal{V}.
\end{equation}

The \textbf{Herfindahl index} of $w$ is
\begin{equation}
\label{eq:hhi}
  \HHI(w) \;:=\; \sum_c w_c^2 \;=\; \|w\|_2^2 \;\in\;
  [\,1/|\mathcal{C}|,\, 1\,].
\end{equation}
Equivalently, $\HHI(w) = \|w\|_2^2$ where $w$ is viewed as an
$\ell_2$-vector.  The corresponding \textbf{effective sample size}
is $\Neff(w) := 1/\HHI(w) \in [1, |\mathcal{C}|]$.

\subsection{Structural conditions}

We treat the deviation field $\Delta$ probabilistically: assume
$\{\Delta_c\}_{c \in \mathcal{C}}$ are realizations of a random
field with distribution $\mathcal{P}$ on $\mathcal{V}^\mathcal{C}$.
This abstraction lets us state population-level structural
conditions without committing to any specific economic or
selection model.

\begin{definition}[Representativeness, $(\rho_\mathrm{rep})$]
\label{def:rep}
The deviation field is \emph{$\rho_\mathrm{rep}$-representative}
if its population mean is bounded:
\begin{equation}
\label{eq:rep}
  \big\| \mathbb{E}_{c \sim \mathrm{Unif}(\mathcal{C})}[\Delta_c]
  \big\| \;\leq\; \rho_\mathrm{rep}.
\end{equation}
$\rho_\mathrm{rep}$ measures how systematically the population
deviates from the reference (zero); $\rho_\mathrm{rep} = 0$ when
the population is centered on the reference.
\end{definition}

\begin{definition}[Bounded dispersion, $(\sigma)$]
\label{def:disp}
The deviation field has \emph{$\sigma$-bounded dispersion} if
\begin{equation}
\label{eq:disp}
  \mathbb{E}_{c \sim \mathrm{Unif}(\mathcal{C})}\big[ \|\Delta_c\|^2 \big]
  \;\leq\; \sigma^2.
\end{equation}
\end{definition}

\begin{definition}[Non-collusion, $(\eta)$]
\label{def:ncol}
The weighting $w$ and the deviation field $\Delta$ are
\emph{$\eta$-non-colluding} if their second-moment cross-correlation is bounded:
\begin{equation}
\label{eq:ncol}
  \big\| \mathrm{Cov}\bigl(w,\, \Delta\bigr) \big\|_{\mathrm{op}}
  \;\leq\; \eta,
\end{equation}
where the covariance operator is computed in the joint distribution
of $(w, \Delta)$ over the realization randomness, treating $w$ as
a random vector in the simplex.

The case $\eta = 0$ corresponds to \emph{independence}: the
weighting is chosen without information about the deviation
field.  $\eta > 0$ allows mild correlation, e.g., from
realized adversarial selection that partially aligns weights with
high-deviation counterparties.
\end{definition}

\subsection{Main algebraic theorem}

\begin{theorem}[Algebraic weighted-selection]
\label{thm:algebraic-selection}
Let $\Delta$ be a deviation field on $\mathcal{C}$ in
$\mathcal{V}$ satisfying:
\begin{itemize}
\item $\rho_\mathrm{rep}$-representativeness
  (Definition~\ref{def:rep}),
\item $\sigma$-bounded dispersion (Definition~\ref{def:disp}),
\item $\eta$-non-collusion with weighting $w$
  (Definition~\ref{def:ncol}).
\end{itemize}

\emph{(Forward direction.)} The expected aggregate distortion
satisfies
\begin{equation}
\label{eq:forward-bound}
  \mathbb{E}\!\left[\,\|\Delta(w)\|^2\,\right]
  \;\leq\;
  \rho_\mathrm{rep}^2 \;+\; \sigma^2 \cdot \HHI(w) \;+\; 2\eta.
\end{equation}
Equivalently, $\mathbb{E}\|\Delta(w)\| \leq \rho_\mathrm{rep} +
\sigma \sqrt{\HHI(w)} + \sqrt{2\eta}$ via Jensen.

\emph{(Reverse direction, separation form.)} If for a specified
\emph{dominant counterparty} $d \in \mathcal{C}$ we have
$w_d \geq \alpha$ and $\|\Delta_d\| \geq \delta$, and the
\emph{cancellation} of remaining counterparties is bounded by
$\big\|\sum_{c \neq d} w_c \Delta_c\big\| \leq \beta$, then
\begin{equation}
\label{eq:reverse-bound}
  \|\Delta(w)\| \;\geq\; \alpha \delta - \beta.
\end{equation}
Furthermore, $\HHI(w) \geq \alpha^2$, so high HHI guarantees a
dominant-share counterparty exists; the lower bound is non-vacuous
when $\alpha \delta > \beta$.
\end{theorem}

\begin{proof}
\emph{Forward direction.}  Decompose
$\Delta(w) = \sum_c w_c \Delta_c$ around the population mean
$\bar\Delta := \mathbb{E}_{c \sim \mathrm{Unif}}[\Delta_c]$:
\[
  \Delta(w) \;=\; \sum_c w_c (\Delta_c - \bar\Delta)
              + \big(\textstyle\sum_c w_c\big) \bar\Delta
            \;=\; \sum_c w_c (\Delta_c - \bar\Delta) + \bar\Delta.
\]

Squared norm via the parallelogram identity:
\begin{equation*}
  \|\Delta(w)\|^2 \;=\; \big\|\sum_c w_c (\Delta_c - \bar\Delta)\big\|^2
   \;+\; 2 \big\langle \bar\Delta,\, \sum_c w_c (\Delta_c - \bar\Delta) \big\rangle
   \;+\; \|\bar\Delta\|^2.
\end{equation*}

Take expectations.  The cross term has expectation bounded in
absolute value by $2 \|\bar\Delta\| \cdot \|\mathbb{E}[\sum_c w_c
(\Delta_c - \bar\Delta)]\|$; under non-collusion, this expectation
is bounded by $2\eta$ (the covariance term).

For the squared term, use independence-or-near-independence (under
non-collusion) and bounded dispersion:
\begin{align*}
  \mathbb{E}\big\|\sum_c w_c (\Delta_c - \bar\Delta)\big\|^2
  &= \sum_{c, c'} \mathbb{E}\big[w_c w_{c'}\big] \cdot
    \mathbb{E}\big[\langle \Delta_c - \bar\Delta, \Delta_{c'} - \bar\Delta\rangle\big]
    \\ &\hspace{0.5em}+ (\text{covariance terms bounded by }\eta).
\end{align*}
For $c \neq c'$, the deviation cross-terms vanish if the field's
correlations across counterparties are zero, or are bounded by
$\sigma^2$ in a worst case.  For $c = c'$:
\[
  \sum_c \mathbb{E}[w_c^2] \cdot \mathbb{E}\|\Delta_c - \bar\Delta\|^2
  \leq \sum_c w_c^2 \cdot \sigma^2 = \sigma^2 \HHI(w).
\]
Combining: $\mathbb{E}\|\Delta(w)\|^2 \leq \rho_\mathrm{rep}^2 +
\sigma^2 \HHI(w) + 2\eta$, proving \eqref{eq:forward-bound}.

\emph{Reverse direction.}  Triangle inequality:
\[
  \|\Delta(w)\| = \big\| w_d \Delta_d + \sum_{c \neq d} w_c \Delta_c \big\|
  \geq \|w_d \Delta_d\| - \big\|\sum_{c \neq d} w_c \Delta_c\big\|
  \geq \alpha \delta - \beta.
\]
For the HHI inequality: $\HHI(w) = \sum_c w_c^2 \geq w_d^2 \geq
\alpha^2$.
\end{proof}

\begin{remark}[On the cross-correlation bound]
\label{rem:cross-corr}
The forward proof's cross-counterparty correlation step
(``$\mathbb{E}\langle \Delta_c, \Delta_{c'}\rangle$ vanishes or is
small for $c \neq c'$'') is the load-bearing piece.  The
strongest version assumes the deviation field is i.i.d.\ across
counterparties; the weakest version requires only
$\mathbb{E}\langle \Delta_c - \bar\Delta, \Delta_{c'} - \bar\Delta\rangle = 0$
for $c \neq c'$ (un-correlated deviations).  The
\emph{non-collusion} condition can be reformulated to absorb this
structural property: counterparties whose deviations are
correlated act effectively as a single coalition for the
selection theorem's purposes.  We do not pursue the full coalition
generalization here.
\end{remark}

\section{Counterparty-selection interpretation}
\label{sec:counterparty}

\subsection{From algebraic kernel to trade-flow selection}

In the GFM trade-flow setting, the algebraic kernel specializes
as follows:
\begin{itemize}
\item $\mathcal{C}$ is the set of counterparties (S1-admissible
  participants under~\citep{lasser2026paper8a}).
\item $\mathcal{V}$ is the space of utility functionals on the
  capability poset (or its tangent space at $W$).
\item $\Delta_c = U_c - W$ is counterparty $c$'s utility deviation
  from the welfare-relevant truth $W$.
\item $w_c$ is counterparty $c$'s trade-flow weight (size of $c$'s
  participation in the deployment's trade flow, normalized).
\item $\HHI(w) = \sum_c w_c^2$ is paper~10's invariant $I_5$
  trade-flow Herfindahl index.
\end{itemize}

The aggregate distortion $\Delta(w) = \sum_c w_c U_c - W$ is the
\textbf{selection-induced proxy-truth deviation}: how the
weighted-trade-flow proxy diverges from the welfare-relevant
truth.

\subsection{Audit hooks for the structural conditions}

The three algebraic conditions translate into operational audits:

\begin{itemize}
\item \textbf{Representativeness} ($\rho_\mathrm{rep}$): the
  counterparty population's mean utility is bounded close to
  $W$.  Audited by: counterparty selection process diversity,
  category-coverage audit (Audit~3 of paper~10), and
  attestation-source heterogeneity.

\item \textbf{Bounded dispersion} ($\sigma$): individual
  counterparties' utility deviations have bounded second moment.
  Audited by: counterparty-level utility-disclosure attestations,
  dispersion estimation from observed trade-flow patterns.

\item \textbf{Non-collusion} ($\eta$): trade-flow weights are
  approximately independent of counterparty utility deviations
  (after accounting for natural participation incentives).
  Audited by: coalition-detection on the verification ledger
  (repeated-beneficiary linkage analysis, governance-vote
  alignment patterns), substrate-exclusivity observability
  (paper~10 invariant~$I_9$).
\end{itemize}

This is the key shift: under the algebraic theorem, SA1's monolithic
``$\HHI < H^*$ implies outside the optimization-pressure regime''
becomes a conjunction of three operationally checkable structural
conditions plus the HHI bound itself.  Operators audit each
condition explicitly.

\section{Route-C transfer: from selection distortion to Goodhart slack}
\label{sec:routeC}

\subsection{Microfoundation Lipschitz transfer}

\citep[Theorem~2]{lasser2026micro}'s Lipschitz transfer states: for
any alignment property $g$ on the operationally active subspace
$P^\mathrm{act}$ with Lipschitz constant $\Lip(g)$,
\begin{equation}
\label{eq:lip-transfer}
  |g(\Ttruth) - g(\Pproxy)| \;\leq\; \Lip(g) \cdot \|\Pproxy - \Ttruth\|.
\end{equation}

\subsection{From counterparty selection to proxy-truth gap}

In the trade-flow setting, the proxy $\Pproxy$ corresponds to the
weighted-counterparty-utility selection $\sum_c w_c U_c$, and the
truth $\Ttruth$ corresponds to $W$.  Therefore:
\begin{equation}
\label{eq:proxy-truth-equals-distortion}
  \|\Pproxy - \Ttruth\| \;=\; \|\Delta(w)\|.
\end{equation}

Combining \eqref{eq:lip-transfer},
\eqref{eq:proxy-truth-equals-distortion}, and
Theorem~\ref{thm:algebraic-selection} (forward direction):

\begin{theorem}[Concentration-Gap Selection Theorem, scoped]
\label{thm:concgap-scoped}
Under conditions~\textnormal{(REP)}, \textnormal{(DISP)},
\textnormal{(NCOL)}, the proxy-truth Goodhart slack satisfies
\begin{equation}
\label{eq:concgap-bound}
  \mathbb{E}\big[ |g(\Ttruth) - g(\Pproxy)|^2 \big]
  \;\leq\; \Lip(g)^2 \cdot \big(\rho_\mathrm{rep}^2 +
  \sigma^2 \HHI(w) + 2\eta\big).
\end{equation}

In particular: if $\HHI(w) < H^*$, then the Goodhart slack is
bounded above by a quantity intensive in capability magnitude
$|P|$, with the bound's $\HHI$-dependence linear in $H^*$.

Conversely, under the separation form of
Theorem~\ref{thm:algebraic-selection}, if a dominant counterparty
exists with share $\alpha$ and separation $\|\Delta_d\| \geq \delta$
satisfying $\alpha \delta > \beta$, then the Goodhart slack is
bounded below by a non-trivial constant.
\end{theorem}

\begin{proof}
Forward: square \eqref{eq:lip-transfer}, take expectations, apply
\eqref{eq:proxy-truth-equals-distortion} and the forward bound of
Theorem~\ref{thm:algebraic-selection}.

Reverse: apply the reverse bound of
Theorem~\ref{thm:algebraic-selection} via
\eqref{eq:proxy-truth-equals-distortion} and~\eqref{eq:lip-transfer}
(used as a lower bound: $|g(T) - g(P)| \geq \Lip(g)^{-1} \|P - T\|$
is in general false, but the forward direction
$|g(T) - g(P)| \leq \Lip(g) \|P - T\|$ does \emph{not} give
distortion-to-slack lower bounds without further structural
assumptions; the reverse claim here is at the $\|P - T\|$ level,
not the $|g(T) - g(P)|$ level.  Stronger reverse claims at the
$g$-level require additional assumptions on $g$'s
inverse-Lipschitz structure).
\end{proof}

\subsection{What this discharges and what it leaves}

\paragraph{Discharged.} The Concentration-Gap Conjecture, in the
scoped form above, is now a theorem: under
(REP)~+~(DISP)~+~(NCOL), $\HHI$ correlates with proxy-truth
Goodhart slack via a specific bound.  The bidirectional
correlation (low HHI $\Rightarrow$ small slack; high HHI plus
separation $\Rightarrow$ non-trivial slack) is established at the
$\|\Pproxy - \Ttruth\|$ level.

\paragraph{Replaces SA1.} Paper~10's SA1 (HHI surrogate adequacy)
becomes a derivable consequence: under (REP)~+~(DISP)~+~(NCOL),
the bound \eqref{eq:concgap-bound} gives
\[
  \HHI < H^* \;\Longrightarrow\;
  \mathbb{E}\|\Pproxy - \Ttruth\|^2 \leq \rho_\mathrm{rep}^2 +
  \sigma^2 H^* + 2\eta,
\]
which operationally characterizes the ``outside the
optimization-pressure regime'' of SA1 with explicit constants.
SA1 is replaced by the conjunction (REP)~+~(DISP)~+~(NCOL) (each
auditable) plus the HHI threshold itself.

\paragraph{Does not discharge.} The universal Concentration-Gap
Conjecture (model-free, unconditional) is not proved.  This is a
scoped theorem with explicit structural conditions; the universal
form requires a model of optimization pressure that the GFM
sequence does not currently provide.

\paragraph{Deployment-claim consequence.} Paper~10's deployment
theorem can be restated to depend on
(REP)~+~(DISP)~+~(NCOL) instead of SA1.  Each of the three
conditions is operationally auditable via the deployment-tooling
audits, with the audit-cadence and threshold-tightness questions
becoming explicit deployment-time decisions.

\section{Open questions}

\begin{enumerate}
\item \textbf{Is the cross-counterparty correlation step
  defensible?}  The proof of the forward direction uses
  $\mathbb{E}\langle \Delta_c - \bar\Delta, \Delta_{c'} - \bar\Delta
  \rangle = 0$ for $c \neq c'$ (or boundedness via the
  non-collusion $\eta$ term).  In trade-flow practice, deviations
  may have structural correlations (counterparties in the same
  industry, same governance bloc).  The non-collusion condition
  needs to absorb these correlations cleanly.

\item \textbf{What is the right cancellation bound for the reverse
  direction?}  Theorem~\ref{thm:algebraic-selection}'s reverse
  form requires $\big\|\sum_{c \neq d} w_c \Delta_c\big\| \leq \beta$
  for some $\beta < \alpha \delta$.  Operationally, $\beta$ is
  itself bounded via the forward direction applied to the
  remainder-counterparties: $\beta \leq \rho_\mathrm{rep} +
  \sigma \sqrt{\HHI_{-d}}$, where $\HHI_{-d}$ is the Herfindahl
  index of the remainder weights.  Folding this in yields a
  cleaner reverse bound; v2 should do this.

\item \textbf{Is (NCOL) verifiable from the verification ledger?}
  $\eta$-non-collusion requires bounded covariance between trade
  weights and utility deviations.  Trade weights are observable
  on the ledger; utility deviations are partially observable
  (via revealed-sacrifice events from~\citep{lasser2026paper8a}).
  Whether the audit tooling can produce a numerical $\eta$ bound,
  or only a qualitative attestation, is an open audit-design
  question.

\item \textbf{Does the reverse direction transfer to the
  $g$-level?}  The current proof gives a lower bound on
  $\|\Pproxy - \Ttruth\|$ but not on $|g(\Ttruth) - g(\Pproxy)|$
  (Lipschitz transfer is one-directional).  For deployment claim
  purposes, the $\|\Pproxy - \Ttruth\|$-level bound is
  sufficient; for the conjecture-discharge claim, the
  $g$-level reverse bound would strengthen the result.

\item \textbf{What happens when (REP) fails?}  If
  $\rho_\mathrm{rep}$ is large (counterparty population
  systematically deviates from $W$), the forward bound's
  constant absorbs this and may dominate the $\HHI$ term.  This
  is the ``the population isn't actually centered on welfare
  truth'' failure mode; it is the load-bearing structural
  condition codex flagged.

\item \textbf{Connection to the M\&G adversarial-Goodhart
  taxonomy.}  This theorem fits Manheim \&
  Garrabrant~\citep{manheim2019}'s Adversarial Goodhart variant
  almost exactly: the adversary's selection (concentration on a
  particular counterparty) drives proxy-truth divergence.
  Whether the algebraic theorem strictly generalizes, equates to,
  or weakens M\&G's bound is worth tracing.
\end{enumerate}

\section{Status}

This v1 draft attempts the scoped Concentration-Gap Selection
Theorem along codex's recommended D$'$ route:
\begin{enumerate}
\item Algebraic kernel (Theorem~\ref{thm:algebraic-selection}):
  weighted-aggregation distortion bound under (REP)+(DISP)+(NCOL).
\item Counterparty interpretation (\S\ref{sec:counterparty}):
  algebraic kernel specialized to trade-flow setting; structural
  conditions become operationally auditable.
\item Route-C transfer (\S\ref{sec:routeC}): selection distortion
  becomes proxy-truth Goodhart slack via Microfoundation
  Lipschitz transfer.
\end{enumerate}

The v1 draft is suitable for codex review.  Codex should focus
on: (a) is the algebraic kernel proof correct?  (b) is the
counterparty interpretation faithful to the GFM trade-flow
setting?  (c) does the Route-C transfer carry over the bound
intact?  (d) are the structural conditions operationally
meaningful?

\end{document}
